\documentclass[12pt,a4paper]{article}
\usepackage{fontspec}
\usepackage{polyglossia}
\setmainlanguage{russian}
\setotherlanguage{english}
\setmainfont{DejaVu Serif}
\setsansfont{DejaVu Sans}
\setmonofont{DejaVu Sans Mono}
\usepackage{amsmath,amssymb}
\usepackage{geometry}
\usepackage{setspace}
\PassOptionsToPackage{hyphens}{url}
\usepackage{hyperref}
\usepackage{graphicx}
\usepackage{microtype}
\usepackage{booktabs}
\usepackage{indentfirst}
\usepackage{enumitem}

\geometry{left=3cm, right=1.5cm, top=2cm, bottom=2cm}
\onehalfspacing
\setlength{\parindent}{1.25cm}

\hypersetup{
  colorlinks=true,
  linkcolor=black,
  urlcolor=blue,
  citecolor=black
}

\begin{document}

\graphicspath{{../grp_forecast/figures/}}
\tolerance=9999
\emergencystretch=4em
\sloppy

%% ---------------------------------------------------------------
%%  Титульный лист
%% ---------------------------------------------------------------
\begin{titlepage}
\begin{center}
{\small Федеральное государственное автономное образовательное учреждение\\
высшего образования\\[4pt]
\textbf{НАЦИОНАЛЬНЫЙ ИССЛЕДОВАТЕЛЬСКИЙ\\
УНИВЕРСИТЕТ\\
«ВЫСШАЯ ШКОЛА ЭКОНОМИКИ»}\\[8pt]
Факультет компьютерных наук\\
Образовательная программа «Анализ больших данных»}

\vspace{3cm}

{\Large\textbf{МАГИСТЕРСКАЯ ДИССЕРТАЦИЯ}}

\vspace{1cm}

{\large\textbf{Прогнозирование валового регионального продукта (ВРП)\\
субъектов Российской Федерации с помощью ансамблевых\\
методов машинного обучения}}

\vspace{3cm}

\begin{flushright}
\begin{tabular}{l}
\textbf{Выполнил:}\\
Семенюченко Артём Геннадьевич\\
студент ФКН НИУ ВШЭ,\\
ОП «Анализ больших данных»\\[12pt]
\textbf{Научный руководитель:}\\
Паточенко Евгений Анатольевич\\
старший преподаватель ФКН НИУ ВШЭ,\\
академический руководитель\\
ОП «Анализ больших данных»
\end{tabular}
\end{flushright}

\vfill
Москва 2026
\end{center}
\end{titlepage}

%% ---------------------------------------------------------------
%%  Аннотация
%% ---------------------------------------------------------------
\newpage
\section*{Аннотация}

Настоящая работа посвящена разработке воспроизводимого подхода
к прогнозированию валового регионального продукта субъектов Российской
Федерации на основе панельных временных данных и ансамблевых методов
машинного обучения. Эмпирическая база~--- региональная панель за 2005--2024
годы; горизонт прогнозирования~--- один год вперёд (h\,=\,1). В роли целевой
переменной рассматривается логарифм реального ВРП. Признаковое пространство
объединяет лаги отклика, производственные и инвестиционные показатели,
индикаторы рынка труда, бюджетные параметры, признаки кризисных периодов
и характеристики региональной специализации.

Сопоставляются три базовые модели~--- наивный прогноз, AR(1) и Ridge
регрессия с фиксированными эффектами~--- и четыре ансамблевых алгоритма:
LightGBM, CatBoost, GPBoost и нейросеть с эмбеддингами регионов (EmbMLP).
Последний из древесных методов~--- GPBoost~--- совмещает градиентный бустинг
со случайными эффектами, явно учитывая панельную структуру данных.
Дополнительно формируется взвешенный ансамбль CatBoost+LightGBM с весами,
оптимизированными post-hoc на отложенных предсказаниях. Качество оценивается
на временно\'{й} валидации по схеме expanding window по метрикам RMSE, MAE
и out-of-sample~$R^2$. Дополнительно проводятся ablation study по блокам
признаков, два robustness check и тест Дибольда--Мариано. Особое внимание
уделено предотвращению leakage: все preprocessing-операции обучаются
исключительно на train-части каждого fold.

\smallskip\noindent
\textbf{Ключевые слова:} валовой региональный продукт, машинное обучение,
панельные временные ряды, ансамблевые модели, градиентный бустинг,
GPBoost, временна\'{я} валидация, SHAP.

%% ---------------------------------------------------------------
%%  Abstract
%% ---------------------------------------------------------------
\newpage
\section*{Abstract}

This thesis develops a reproducible pipeline for forecasting the Gross Regional
Product (GRP) of Russian Federation subjects using panel time-series data and
ensemble machine learning methods. The empirical base is a regional panel
covering 2005--2024; the forecast horizon is one year ahead ($h=1$). The target
variable is defined as the logarithm of real GRP. The feature space combines
lagged values of the target, industrial and investment indicators, labour market
statistics, fiscal variables, crisis dummies, and regional specialisation
characteristics.

The study compares three baseline models~--- naive forecast, AR(1), and Ridge
regression with fixed effects~--- and four ensemble algorithms: LightGBM,
CatBoost, GPBoost, and an embedding-based neural network (EmbMLP).
GPBoost combines gradient boosting with Gaussian random effects, explicitly
addressing panel heterogeneity without expanding the feature space.
Additionally, a weighted CatBoost+LightGBM ensemble is constructed with
post-hoc optimised weights ($w^*=0{.}25$ for CatBoost, $0{.}75$ for LightGBM).
Model quality is assessed via expanding-window time-series cross-validation
using RMSE, MAE, and out-of-sample~$R^2$. The experimental design includes an
ablation study across feature groups, two robustness checks, and a
Diebold--Mariano test. Particular attention is paid to leakage prevention: all
preprocessing operations are fitted exclusively on the training portion of each
fold.

\smallskip\noindent
\textbf{Keywords:} gross regional product, machine learning, panel time series,
ensemble methods, gradient boosting, GPBoost, time-series validation, SHAP.

%% ---------------------------------------------------------------
%%  Содержание
%% ---------------------------------------------------------------
\newpage
\tableofcontents

%% ---------------------------------------------------------------
%%  1. ВВЕДЕНИЕ
%% ---------------------------------------------------------------
\newpage
\section{Введение}

\subsection{Актуальность исследования}

Валовой региональный продукт занимает центральное место среди показателей
экономического развития субъектов Российской Федерации. По экономическому
содержанию ВРП близок к ВВП, но рассчитывается для отдельной территории,
что позволяет перейти от агрегированного описания к анализу пространственного
распределения выпуска и неоднородной реакции регионов на экономические шоки.

Российская экономика характеризуется выраженной территориальной
неоднородностью, обусловленной историей экономического становления
страны~--- централизованностью вокруг крупных городов и большим количеством
моногородов~--- а также особенностями ландшафта и климата (регионы-курорты,
добывающие регионы). В результате в одной стране сосуществуют сырьевые,
индустриальные, аграрные и бюджетно-зависимые регионы, каждый из которых
обладает собственным механизмом формирования валового продукта.
Прикладную значимость исследования усиливает лаг публикации официальной
статистики: по данным Росстата, объём ВРП по сумме субъектов за 2024 год
в текущих ценах составил 186\,545,1 млрд руб., индекс физического
объёма~--- 104,4\%~\cite{rosstat2024}. Эти оценки появляются после завершения
отчётного периода, тогда как региональный мониторинг и бюджетное планирование
требуют более оперативной картины динамики выпуска.

Сама задача относится к классу регрессии на панельных временных данных.
Наблюдением выступает пара «регион~--- год», целевой переменной~--- будущее
значение ВРП, признаками~--- лаговая информация и социально-экономические
показатели региона. Такая структура отличается от классического одномерного
временного ряда: модель должна одновременно учитывать временну\'{ю} зависимость,
межрегиональную неоднородность, различие масштабов субъектов и возможные
структурные разрывы. Дополнительная сложность состоит в риске утечки
информации из будущего: применение стандартных процедур кросс-валидации
без учёта временного порядка систематически завышает оценки качества.
Cerqua, Letta и Pinto~\cite{cerqua2025} показывают, что leakage в панельных
данных имеет несколько источников, не сводимых к простому «подглядыванию»
в тестовую выборку.

Современная литература по macro-ML подчёркивает, что применение машинного
обучения оправдано лишь при строгой экспериментальной процедуре~\cite{chi2025}.
Соответственно, ансамблевые методы рассматриваются в данной работе не как
заведомо более точные, а как гипотеза, подлежащая проверке относительно
сильных базовых моделей.

\subsection{Объект и предмет исследования}

\textit{Объектом} исследования выступают панельные временные ряды
социально-экономических показателей субъектов Российской Федерации.
\textit{Предметом}~--- методы прогнозирования ВРП на этих данных,
охватывающие как эконометрические базовые модели, так и ансамблевые
алгоритмы машинного обучения.

\subsection{Цель и задачи исследования}

\textit{Цель работы}~--- разработка и оценка воспроизводимого пайплайна
прогнозирования ВРП субъектов Российской Федерации на панельных временных
данных с использованием эконометрических базовых моделей и ансамблевых
методов машинного обучения.

Для достижения цели решаются следующие задачи:
\begin{enumerate}
  \item Провести обзор литературы по макроэкономическому и региональному
        прогнозированию, включая современные ML-подходы.
  \item Сформировать панельную базу данных за 2005--2024 годы, проверить
        стационарность и определить целевую переменную.
  \item Реализовать предобработку данных и leakage-free процедуру
        формирования признаков.
  \item Построить базовые модели (наивный прогноз, AR(1), Ridge с FE)
        и ансамблевые алгоритмы (LightGBM, CatBoost, GPBoost).
  \item Реализовать временну\'{ю} валидацию по схеме expanding window
        и подбор гиперпараметров строго внутри train-части каждого fold.
  \item Сравнить модели по RMSE, MAE и out-of-sample~$R^2$.
  \item Провести ablation study по блокам признаков и robustness checks.
  \item Построить SHAP-интерпретацию лучшей модели.
\end{enumerate}

\subsection{Данные и инструментарий}

Информационная база формируется на основе данных Росстата, ЕМИСС,
Банка России, Минэкономразвития и Федерального казначейства за 2005--2024
годы. В признаковое пространство включены лаги ВРП, индекс промышленного
производства, показатели добычи и обработки, сельское хозяйство,
строительство, транспорт, торговля, инвестиции в основной капитал, занятость,
безработица, заработная плата, доходы населения, доходы и расходы
региональных бюджетов, трансферты, федеральный округ и dummy-переменные
кризисных периодов.

Используемые библиотеки: \texttt{pandas~2.2}, \texttt{numpy~1.26}~--- работа
с данными; \texttt{scikit-learn~1.4}~--- линейные модели, валидация,
масштабирование; \texttt{lightgbm~4.3}, \texttt{catboost~1.2}~---
бустинговые алгоритмы; \texttt{gpboost~1.5}~--- модель со смешанными
эффектами; \texttt{linearmodels~6.0}~--- панельные регрессии с
фиксированными эффектами; \texttt{optuna~3.6}~--- подбор гиперпараметров;
\texttt{shap~0.45}~--- интерпретация; \texttt{matplotlib~3.8},
\texttt{seaborn~0.13}~--- визуализация. Реализация выполнена на Python~3.11.
Воспроизводимость обеспечивается фиксацией \texttt{random\_state=42} во всех
стохастических процедурах и сохранением конфигураций в YAML-файлах.

\subsection{Практическая значимость}

Результаты могут использоваться для мониторинга региональной экономической
динамики, предварительной оценки ВРП до публикации официальной статистики
и построения сценарных прогнозов. SHAP-разложения позволяют сопоставить
выявленные факторы с экономической логикой регионов, что делает модель
пригодной для содержательного аналитического использования.

\subsection{Научная новизна}

В существующих российских исследованиях акцент делается либо на
детерминантах ВРП~\cite{badykova2026}, либо на сценарном прогнозе для
отдельного региона~\cite{malkina2025}. Новизна настоящей работы~--- в
формировании единого экспериментального фреймворка: прогнозирование ВРП
как регрессия на панельных временных данных с временно\'{й} валидацией,
явным контролем leakage, ablation study, robustness checks и
SHAP-интерпретацией на единой панели субъектов. Дополнительно в сравнение
включён GPBoost~--- алгоритм, совмещающий градиентный бустинг со случайными
эффектами и ранее не применявшийся к задаче прогнозирования ВРП российских
регионов.

%% ---------------------------------------------------------------
%%  2. ОБЗОР ЛИТЕРАТУРЫ
%% ---------------------------------------------------------------
\section{Обзор литературы}

Задача прогнозирования ВРП пересекается с несколькими исследовательскими
направлениями: макроэкономическое прогнозирование, машинное обучение на
панельных данных и интерпретируемые модели. Ниже детально разбираются
четыре работы, наиболее значимых для методологии настоящего исследования,
а затем кратко резюмируется сопутствующая литература.

\subsection{Chi et al.\ (2025). Macroeconomic Forecasting and
Machine Learning~\cite{chi2025}}

\textit{Постановка.} В рамках данной работы акцент смещается с точечного
прогнозирования среднего на условное распределение макроэкономических
результатов с помощью машинного обучения. Рассматривается задача
прогнозирования уровня безработицы с длительным горизонтом. Методологически
статья объединяет три элемента современной прогнозной практики: работу с
высокоразмерными данными, регуляризацию для контроля сложности модели и
строгую out-of-sample валидацию.

\textit{Метод и ключевые результаты.} В статье используются глубокие
нейронные сети для прогнозирования квантилей распределения. В качестве
функции потерь применяется квантильная регрессия (так называемая, pinball loss). Авторы
применяют схему расширяющееся окна (expanding-window), в рамках которой модель обучается только на прошлых
наблюдениях, гиперпараметры подбираются на валидационном периоде, а качество в свою очередь
проверяется на тестовом периоде. В конечном итоге, как результат, регуляризация играет
решающую роль. Без сжатия высокоразмерная модель извлекает шум, вместо ожидаемого
устойчивого макроэкономического сигнала. Нелинейные преобразования дают лишь
ограниченный прирост точности~--- основную роль играет
процедура отбора сложности самой модели и корректная out-of-sample валидация.

\textit{Значимость для настоящей работы.} Статья, своего рода, становится референсом в вопросе того как приняется ML в макроэкономике. Сложные модели сравниваются
с базовыми только во вневыборочной процедуре. Вывод об ограниченной роли
нелинейности предостерегает от автоматического ожидания превосходства
ансамблевых методов над линейными, что очень схоже с раскладом в рамках текущего исследования

\subsection{Badykova (2026). Детерминанты ВРП в России:
подход машинного обучения~\cite{badykova2026}}

\textit{Постановка.} Анализ автором детерминанты ВРП на душу населения
по 85 регионам России за 2013--2023 годы. Основной целью является интерпретация, а не  прогнозирование.

\textit{Метод и результаты.} Основным алгоритмом выступает LightGBM.
Интерпретация построена на SHAP: глобальный анализ показывает важность
признаков по всей выборке, локальный~--- объясняет отклонения конкретных
регионов. Наиболее значимыми предикторами оказались инвестиции в основной
капитал и занятость в обрабатывающей промышленности.

\textit{Отличие от настоящей работы.} Badykova решает задачу объяснения,
а не прогнозирования. Настоящая работа переориентирована на вневыборочный
прогноз с временно\'{й} валидацией, tuning только внутри train, ablation
study и явным контролем leakage. Кроме того, в сравнение добавлен GPBoost.

\subsection{Lashina \& Grishunin (2023). Сравнение моделей
прогнозирования ВВП России~\cite{lashina2023}}

\textit{Постановка.} Авторы сопоставляют прогностическую точность четырёх
классов моделей применительно к ВВП России в условиях кризисных периодов:
ARIMA, VAR, SVR и CatBoost. Исследование охватывает кризисы 2014--2015
и 2020 годов.

\textit{Ключевые результаты.} CatBoost демонстрирует наилучшее качество
в кризисные периоды, превосходя ARIMA и VAR по RMSE и MAPE. В стабильные
периоды разница сокращается. Работа является единственной в рассмотренной
литературе, сравнивающей CatBoost с классическими эконометрическими
моделями на данных ВВП России и охватывающей кризисные периоды.

\textit{Значимость.} Результаты обосновывают включение CatBoost в сравнение
и мотивируют отдельный анализ качества моделей в кризисные годы (robustness
check).

\subsection{Sigrist (2022). Gaussian Process Boosting~\cite{sigrist2022}}

\textit{Постановка.} Работа предлагает алгоритм GPBoost, объединяющий
градиентный бустинг с гауссовыми процессами и моделями смешанных эффектов.
Мотивация: стандартный бустинг моделирует только условное среднее
$\mathbb{E}[y \mid X]$, тогда как в панельных данных остатки систематически
коррелированы внутри групп (регионов), что снижает эффективность оценки.

\textit{Модель.} Совместная спецификация:
\begin{equation}
  y_{i,t} = F(X_{i,t}) + b_i + \varepsilon_{i,t},
  \quad b_i \sim \mathcal{N}(0,\sigma_b^2),
  \quad \varepsilon_{i,t} \sim \mathcal{N}(0,\sigma^2),
\end{equation}
где $F(\cdot) = \sum_m \eta T_m(\cdot)$~--- аддитивная функция деревьев,
$b_i$~--- случайный эффект региона~$i$.

\textit{Значимость для настоящей работы.} GPBoost включён как алгоритм,
специально предназначенный для панельной структуры. Это позволяет проверить,
даёт ли явное моделирование региональной гетерогенности прирост качества
сверх стандартного бустинга.

\subsection{Сопутствующая литература}

\textit{Наукастинг.} Цукарев и соавторы~\cite{tsukarev2025} сравнивают
методы bridge equations и MIDAS с ML-алгоритмами на данных Армении
и Беларуси, показывая преимущество бустинга на агрегированных годовых данных.
Макеева и Станкевич~\cite{makeeva2022} применяют MIDAS и MFBVAR для
оперативной оценки ВВП России.

\textit{Региональные исследования.} Малкина, Сочков и Капустина~\cite{malkina2025}
разрабатывают нейросетевые сценарные прогнозы ВРП Нижегородской области.
Банк России~\cite{semiturkin2022} показывает, что целесообразность ML
следует проверять отдельно для каждой группы регионов.

\textit{Интерпретируемость и ML-дизайн.} Lundberg и Lee~\cite{lundberg2017}
разработали SHAP-разложение; в настоящей работе SHAP применяется для
проверки экономической согласованности модели, а не как инструмент
причинной идентификации. Смирнов~\cite{smirnov2025} предупреждает, что
прирост ML относительно классических моделей на традиционных данных
нередко умеренный~--- этот тезис усиливает необходимость строгого сравнения
с наивным прогнозом и AR(1).

%% ---------------------------------------------------------------
%%  3. МЕТОДОЛОГИЯ
%% ---------------------------------------------------------------
\section{Методология исследования}

\subsection{Формальная постановка задачи}

Пусть $i \in \{1,\ldots,N\}$~--- субъект Федерации, $t \in \{1,\ldots,T\}$~---
год, $h=1$~--- горизонт прогноза. Задача состоит в нахождении функции
$f_\theta$ такой, что
\begin{equation}
  \hat{y}_{i,t+1} = f_\theta(X^{\text{available}}_{i,t}),
\end{equation}
где $y_{i,t} = \log(\mathrm{GRP}^{\text{real}}_{i,t})$~--- логарифм реального
ВРП, $X^{\text{available}}_{i,t}$~--- вектор признаков, содержащий только
информацию, доступную к моменту~$t$. Множество данных:
\begin{equation}
  \mathcal{D} = \{(X_{i,t}, y_{i,t+1}) : i = 1,\ldots,N;\; t = p,\ldots,T-1\},
\end{equation}
где $p$~--- максимальный лаг признаков.

\subsection{Проверка стационарности}

Перед построением моделей проводится проверка стационарности. Для каждого
региона применяется расширенный тест Дики~--~Фуллера (ADF) на логарифм ВРП
и его первую разность. На уровне панели используется тест Im~--~Pesaran~--~Shin
(IPS), допускающий неодинаковые авторегрессионные коэффициенты по регионам.
Гипотеза о нестационарности уровня ВРП ожидаемо не отвергается; первые
разности предположительно стационарны. Эти результаты обосновывают
включение лагов ВРП как признаков.

\subsection{Признаковое пространство}

Признаковая система строится как экономически мотивированная структура:
\begin{equation}
  X_{i,t} = \bigl[X^{\text{lag}}_{i,t},\; X^{\text{prod}}_{i,t},\;
  X^{\text{inv}}_{i,t},\; X^{\text{lab}}_{i,t},\; X^{\text{bud}}_{i,t},\;
  X^{\text{cat}}_{i,t},\; X^{\text{shock}}_t\bigr].
\end{equation}

\textit{Инерционный блок} объединяет лаги ВРП $y_{i,t-1}$, $y_{i,t-2}$
и скользящие средние. \textit{Производственный}~--- показатели добычи,
обработки, промышленности, сельского хозяйства, строительства, транспорта
и торговли. \textit{Инвестиционный}~--- инвестиции в основной капитал,
инвестиции на душу населения, отношение инвестиций к ВРП.
\textit{Социально-трудовой}~--- занятость, безработица, заработная плата,
доходы населения. \textit{Бюджетный}~--- доходы и расходы бюджета,
трансферты, бюджетные инвестиции. \textit{Категориальный блок}~---
федеральный округ и индикаторы специализации (сырьевая, индустриальная,
аграрная, сервисная, бюджетно-зависимая). \textit{Блок шоков}~--- dummy
кризисных периодов 2008--2009, 2014--2015, 2020 и 2022--2023 годов.

\subsection{Предотвращение leakage и предобработка}

Применительно к панельным данным leakage имеет несколько
источников~\cite{cerqua2025}. Базовое правило:
\begin{equation}
  X^{\text{available}}_{i,t} = \{x^{(j)}_{i,s} : s \le t - \ell_j,\;
  j = 1,\ldots,p\},
\end{equation}
где $\ell_j \ge 1$~--- лаг доступности признака. Все операции нормализации
и импутации обучаются исключительно на $\mathrm{Train}_k$ и применяются
к $\mathrm{Test}_k$. Денежные показатели дефлируются индексами
потребительских цен. Для масштабирования применяется \texttt{RobustScaler}:
\begin{equation}
  x_{\text{scaled}} = \frac{x - Q_2}{Q_3 - Q_1},
\end{equation}
где $Q_1, Q_2, Q_3$~--- квартили обучающей выборки. Выбор \texttt{RobustScaler}
перед \texttt{StandardScaler} обусловлен структурными выбросами
(ХМАО, ЯНАО, Москва), не являющимися ошибками измерения, но способными
смещать оценку дисперсии.

\subsection{Базовые модели}

В работе используются три базовые модели, представляющие нарастающую
сложность спецификаций.

\textit{Наивный прогноз} воспроизводит последнее известное значение:
$\hat{y}_{i,t+1} = y_{i,t}$.

\textit{AR(1) по регионам} оценивается отдельно для каждого субъекта:
$y_{i,t+1} = \alpha_i + \rho_i y_{i,t} + \varepsilon_{i,t+1}$.

\textit{Ridge с фиксированными эффектами} объединяет индивидуальные
константы региона и временны\'{х} периодов с L2-регуляризацией полного
набора признаков:
\begin{equation}
  \hat{\beta}^{\text{ridge}} = \arg\min_{\beta}
  \Bigl[\sum_{i,t}(y_{i,t+1} - \alpha_i - \delta_t - X^\top_{i,t}\beta)^2
  + \lambda\|\beta\|^2_2\Bigr].
\end{equation}

\subsection{Ансамблевые модели}

В ходе исследования сравниваются три ансамблевых алгоритма:

\textit{LightGBM}~\cite{ke2017} реализует градиентный бустинг на
деревьях с leaf-wise стратегией роста, разбиением признаков
(EFB) и гистограммным методом. Выбран как основной вариант бустинга

\textit{CatBoost}~\cite{dorogush2018} устраняет target leakage при обработке категориальных признаков, а также использует симметричные деревья. По результатам Lashina \& Grishunin~\cite{lashina2023}
демонстрирует наилучшее качество на российских данных в кризисные периоды.

\textit{GPBoost}~\cite{sigrist2022} совмещает бустинг со случайными
эффектами регионов:
\begin{equation}
  y_{i,t} = F(X_{i,t}) + b_i + \varepsilon_{i,t},
  \quad b_i \sim \mathcal{N}(0,\sigma_b^2),
  \quad \varepsilon_{i,t} \sim \mathcal{N}(0,\sigma^2).
\end{equation}
Выбран на основе ранее проделанных исследований коллег в качестве альтернативы добавлению бинарных переменных регионов. Разница между GPBoost и их (dummy-переменных) включением заключется в тотм, что бустинг оценивает
случайные эффекты ковариационно, не увеличивая размерность признакового
пространства и позволяет при этом строить прогнозы для регионов с короткой историей
наблюдений.

\textit{Взвешенный ансамбль} Комбинация из CatBoost \& LightGBM предсказаний
через взвешенную сумму следующим образом:
\begin{equation}
  \hat{y}^{\mathrm{ens}}_{i,t} = w\,\hat{y}^{\mathrm{CB}}_{i,t}
  + (1-w)\,\hat{y}^{\mathrm{LGB}}_{i,t}.
\end{equation}
Вес~$w$ задаётся post-hoc посредством перебора
$w \in \{0{,}00,\,0{,}05,\,\ldots,\,1{,}00\}$: для каждого значения
вычисляется pooled RMSE на объединённой hold-out-выборке девяти fold'ов:
\begin{equation}
  \mathrm{RMSE}(w) = \sqrt{\frac{1}{n}\sum_{i,t}\bigl(y_{i,t}
  - w\,\hat{y}^{\mathrm{CB}}_{i,t}
  - (1-w)\,\hat{y}^{\mathrm{LGB}}_{i,t}\bigr)^2}.
\end{equation}
Минимум достигается при $w^* = 0{,}25$ ($\mathrm{RMSE}(w^*) = 0{,}0382$),
что свидетельствует о более высоком предсказательном вкладе LightGBM.

\subsection{Временна\'{я} валидация и подбор гиперпараметров}

Основная схема~--- расширяющееся окно:
\begin{equation}
  \mathrm{Train}_k = \{(i,t) : t \le \tau_k\},\quad
  \mathrm{Test}_k = \{(i,t) : t = \tau_k + 1\}.
\end{equation}
Начальное окно обучения~--- не менее десяти лет (2005--2014),
$\tau_k$ последовательно увеличивается. Итоговый размер~--- восемь--десять
fold'ов.

Подбор гиперпараметров выполняется через дополнительное внутреннее
разбиение $\mathrm{Train}_k = \mathrm{Train}^{\text{inner}}_k \cup
\mathrm{Val}^{\text{inner}}_k$ с условием
$\max(\mathrm{Train}^{\text{inner}}_k) < \min(\mathrm{Val}^{\text{inner}}_k)$.
Параметры выбираются через Optuna (TPE sampler), после чего модель
переобучается на полном $\mathrm{Train}_k$.

Взвешенный ансамбль обучается в два этапа.
\begin{enumerate}[leftmargin=*,noitemsep]
  \item \textit{Независимая настройка и обучение подмоделей.}
        CatBoost и LightGBM настраиваются отдельно в рамках той же
        expanding-window схемы: 60~trials Optuna на каждый fold,
        внутреннее валидационное окно~--- последние два года $\mathrm{Train}_k$.
        Предсказания каждой подмодели на тестовой части fold'а сохраняются
        в parquet-файлах.
  \item \textit{Post-hoc оптимизация веса.}
        После завершения CV по сохранённым предсказаниям всех девяти fold'ов
        проводится перебор $w \in \{0{,}00,\,0{,}05,\,\ldots,\,1{,}00\}$
        с минимизацией pooled RMSE (см.~формулу~(5)).
        Вес $w^* = 0{,}25$ определяется исключительно по hold-out-данным~---
        целевая переменная обучающей выборки не задействуется, что полностью
        исключает утечку.
\end{enumerate}

\subsection{Метрики качества}

Основные метрики для сравнения моделей:
\begin{align}
  \mathrm{RMSE} &= \sqrt{\frac{1}{n}\sum_j(y_j - \hat{y}_j)^2}, &
  \mathrm{MAE} &= \frac{1}{n}\sum_j|y_j - \hat{y}_j|, \\[4pt]
  R^2_{\mathrm{OOS}} &= 1 - \frac{\sum_j(y_j-\hat{y}_j)^2}
  {\sum_j(y_j-\hat{y}^{\text{bench}}_j)^2},
\end{align}
где в знаменателе $R^2_{\mathrm{OOS}}$ в качестве benchmark используется
наивный прогноз. Положительный $R^2_{\mathrm{OOS}}$ означает, что модель
превосходит наивный прогноз за пределами обучающей выборки.

%% ---------------------------------------------------------------
%%  4. РЕЗУЛЬТАТЫ ЭКСПЕРИМЕНТОВ
%% ---------------------------------------------------------------
\section{Результаты экспериментов}

\subsection{Сравнение моделей на скользящей кросс-валидации}

Кросс-валидация выполнена на расширяющемся окне из девяти fold'ов
(тестовые годы 2015--2023) на панели 85 субъектов Российской Федерации
за период 2005--2023. Целевая переменная~--- темп прироста логарифма
реального ВРП: $\Delta\log\mathrm{VRP}_{i,t} = \log\mathrm{VRP}_{i,t} - \log\mathrm{VRP}_{i,t-1}$.
Бенчмарком служит наивный прогноз (среднее темпа роста по обучающим годам).

\begin{center}
\small
\begin{tabular}{lrrr}
\toprule
Модель & RMSE & MAE & $R^2_{\mathrm{OOS}}$ \\
\midrule
Ансамбль CB+LGBM$^*$ & \textbf{0.0382$\pm$0.013} & \textbf{0.0288$\pm$0.011} & \textbf{0.144$\pm$0.33} \\
LightGBM        & 0.0388$\pm$0.012 & 0.0295$\pm$0.011 & 0.122$\pm$0.26 \\
CatBoost        & 0.0408$\pm$0.015 & 0.0311$\pm$0.014 & 0.036$\pm$0.39 \\
Наивный прогноз & 0.0419$\pm$0.012 & 0.0331$\pm$0.010 & 0.000 \\
AR(1)           & 0.0487$\pm$0.022 & 0.0336$\pm$0.012 & $-$0.404$\pm$0.89 \\
EmbMLP          & 0.0491$\pm$0.020 & 0.0378$\pm$0.017 & $-$0.439$\pm$0.92 \\
Ridge с FE      & 0.0517$\pm$0.026 & 0.0419$\pm$0.027 & $-$0.618$\pm$1.09 \\
GPBoost         & 0.4882$\pm$0.596 & 0.4285$\pm$0.557 & $-$280.9$\pm$515.5 \\
\bottomrule
\multicolumn{4}{l}{\scriptsize $^*$Веса: CB~$\times$~0{,}25, LGB~$\times$~0{,}75 (оптимум grid search, $w\!\in\![0{,}1]$, шаг~0{,}05).}
\end{tabular}
\end{center}

Лучший результат показал ансамбль CatBoost+LightGBM с оптимальными
весами 0{,}25/0{,}75 (соотношение найдено перебором по отложенной выборке):
RMSE~=~0{,}0382, снижение относительно наивного прогноза составило 8{,}8\%;
$R^2_{\mathrm{OOS}} = 0{,}144 > 0$ подтверждает устойчивое превосходство над
бенчмарком вне обучающей выборки. LightGBM (RMSE~=~0{,}0388) занял второе
место с отрывом 0{,}0006 от ансамбля; CatBoost~--- третий
(RMSE~=~0{,}0408, улучшение над наивным 2{,}6\%). AR(1), EmbMLP и Ridge с FE
уступают наивному прогнозу; GPBoost показал нестабильные результаты
из-за технических ограничений при OOS-прогнозировании (см.~раздел~5).

\begin{figure}[htbp]
  \centering
  \includegraphics[width=0.6\textwidth]{metrics_table}
  \caption{Сводная таблица метрик по всем моделям.}
  \label{fig:metrics_table}
\end{figure}

\begin{figure}[htbp]
  \centering
  \includegraphics[width=\textwidth]{rmse_and_growth_combined}
  \caption{Фактический рост ВРП в обучающем периоде 2006--2014 (левая часть)
           и RMSE прогнозов по тестовым годам 2015--2023 (правая часть) для
           всех моделей, включая оптимальный ансамбль
           CB$\times$0{,}25+LGB$\times$0{,}75.}
  \label{fig:combined}
\end{figure}

\begin{figure}[htbp]
  \centering
  \includegraphics[width=0.65\textwidth]{ensemble_weight_search}
  \caption{Зависимость pooled RMSE ансамбля от веса CatBoost~$w$
           ($1-w$~--- доля LightGBM). Минимум при $w^*=0{,}25$;
           при $w=1{,}0$ (чистый CatBoost) RMSE\,=\,0{,}0408,
           при $w=0{,}0$ (чистый LightGBM) RMSE\,=\,0{,}0388.}
  \label{fig:weight_search}
\end{figure}

\subsection{Анализ влияния компонентов (Ablation study)}

Подразумевает под собой последовательное добавление блоков признаков к LightGBM~---
лучшей одиночной модели:

\begin{itemize}
  \item \textbf{A0}~--- только лаги ВРП (LAG\_BLOCK);
  \item \textbf{A1}~--- A0 + производственный блок;
  \item \textbf{A2}~--- A1 + инвестиционный, социально-трудовой и бюджетный блоки;
  \item \textbf{A3}~--- полный набор (A2 + макро, цены, категориальные признаки,
                         кризисные dummy).
\end{itemize}

Во всех представленных кейсах RMSE практически схоже с наивным прогнозом ($\approx$~0{,}043--0{,}044).
Наибольший прирост происходит в рамках перехода A0~$\to$~A1: за добавления производственного блока, снижается RMSE приблизительно с~0{,}044 до~0{,}043. Сценарии A2 и A3 практически не отличимы от наивного прогноза, что объясняется большой долей пропусков в ряде бюджетных показателей. Данный результат идет в конфронтацию с поведением CatBoost,
для которого полная спецификация A3 обеспечивает более выраженное снижение ошибки, по причине того, что он лучше использует категориальные и макро-переменные.

\begin{figure}[htbp]
  \centering
  \includegraphics[width=0.75\textwidth]{ablation_study}
  \caption{Ablation study: прирост точности прогноза при последовательном
           добавлении блоков признаков к LightGBM.}
  \label{fig:ablation}
\end{figure}

\subsection{Анализ устойчивости}

Проведены два robustness check.

\textit{RC1~--- исключение регионов-выбросов.} Из обучающей и тестовой
выборок удалены четыре региона, доминирующих по масштабу ВРП: Москва,
Ханты-Мансийский АО, Ямало-Ненецкий АО и Сахалинская область.
Ансамбль CB$\times$0{,}25+LGB$\times$0{,}75 сохранил лидерство;
среди одиночных древесных моделей LightGBM опередил CatBoost,
что подтверждает: результаты не обусловлены несколькими структурными выбросами.

\textit{RC2~--- кризисные против нормальных периодов.} Fold'ы разделены
на кризисные (2008--2009, 2014--2015, 2020, 2022--2023) и некризисные годы.
Ансамбль остаётся наилучшей моделью в обеих подгруппах, подтверждая тезис
Lashina \& Grishunin~\cite{lashina2023} о преимуществе бустинга в условиях
структурных разрывов.

\begin{figure}[htbp]
  \centering
  \includegraphics[width=0.48\textwidth]{rc1_outliers}
  \hfill
  \includegraphics[width=0.48\textwidth]{rc2_crisis}
  \caption{RC1: сравнение метрик с~выбросами и~без~них (слева).
           RC2: RMSE в~кризисные и~некризисные периоды (справа).}
  \label{fig:robustness}
\end{figure}

\subsection{SHAP-интерпретация}

Для LightGBM (лучшей одиночной модели) построено SHAP-разложение:
\begin{equation}
  \hat{y}_{i,t} = \phi_0 + \sum_{j=1}^{p} \phi_{j,i,t},
\end{equation}
где $\phi_{j,i,t}$~--- вклад признака $j$ в прогноз для региона $i$ в
году $t$. Глобальная значимость: $I_j = \frac{1}{n}\sum_{i,t}|\phi_{j,i,t}|$.

Пять признаков с наибольшим средним $|\phi_j|$:
\begin{enumerate}
  \item \texttt{usd\_rub\_volatility\_lag1}~--- внутригодовая волатильность
        курса USD/RUB; вклад превышает следующий признак примерно вдвое;
  \item \texttt{min\_food\_basket\_cost\_lag1}~--- стоимость минимального
        набора продуктов (потребительский ценовой индикатор);
  \item \texttt{policy\_rate\_lag1}~--- номинальная ключевая ставка ЦБ РФ;
  \item \texttt{policy\_rate\_real\_lag1}~--- реальная ключевая ставка;
  \item \texttt{paid\_services\_to\_grp\_lag1}~--- доля платных услуг в ВРП.
\end{enumerate}
Валютный блок (PRICE\_BLOCK) и макро-монетарный блок (MACRO\_BLOCK) совместно
доминируют над отраслевыми и инвестиционными показателями. Высокая волатильность
USD/RUB в предшествующем году сигнализирует о монетарной неопределённости,
подавляющей региональные инвестиции и, как следствие, рост ВРП. Ставка ЦБ
(номинальная и реальная) отражает канал кредитных условий. Региональный тип
(CATEGORICAL\_BLOCK) имеет нулевой вклад во всех типах регионов, тогда как
лаговые компоненты ВРП (LAG\_BLOCK) дают лишь~0{,}09--0{,}10~---
подтверждая низкую автокорреляцию годовых темпов роста ВРП.

\begin{figure}[htbp]
  \centering
  \includegraphics[width=0.48\textwidth]{shap_global}
  \hfill
  \includegraphics[width=0.48\textwidth]{shap_beeswarm}
  \caption{Глобальный SHAP-анализ LightGBM: среднее $|\phi_j|$
           по признакам (слева) и распределение SHAP-значений
           по наблюдениям (справа).}
  \label{fig:shap_global}
\end{figure}

\begin{figure}[htbp]
  \centering
  \includegraphics[width=0.7\textwidth]{shap_group_heatmap}
  \caption{Групповой SHAP: вклад блоков признаков в разрезе типов регионов.}
  \label{fig:shap_group}
\end{figure}

Локальный анализ выполнен для пяти репрезентативных субъектов:
Санкт-Петербург, ЯНАО, Тамбовская область, Чувашия и Ярославская область.
Они представляют различные экономические профили~--- от ресурсных до
бюджетно-зависимых~--- и позволяют верифицировать направление SHAP-эффектов.

\begin{figure}[htbp]
  \centering
  \includegraphics[width=0.48\textwidth]{shap_local_spb}
  \hfill
  \includegraphics[width=0.48\textwidth]{shap_local_yanao}
  \caption{Локальные SHAP-waterfall: Санкт-Петербург (слева) и ЯНАО (справа).}
  \label{fig:shap_local_1}
\end{figure}

\begin{figure}[htbp]
  \centering
  \includegraphics[width=0.48\textwidth]{shap_local_tambov}
  \hfill
  \includegraphics[width=0.48\textwidth]{shap_local_chuvashiya}\\[6pt]
  \includegraphics[width=0.48\textwidth]{shap_local_yaroslavl}
  \caption{Локальные SHAP-waterfall: Тамбовская обл. (слева), Чувашия (справа),
           Ярославская обл. (снизу).}
  \label{fig:shap_local_2}
\end{figure}

\subsection{Анализ ошибок}

Ошибки ансамбля CB$\times$0{,}25+LGB$\times$0{,}75 проанализированы в двух разрезах.

По регионам ($\mathrm{MAE}_i = \frac{1}{T}\sum_t|y_{i,t}-\hat{y}_{i,t}|$):
наибольшие ошибки характерны для регионов с нестационарной структурой
экономики и высокой волатильностью ВРП. Детальный рейтинг приведён на
рисунке~\ref{fig:error_region}.

По годам ($\mathrm{MAE}_t = \frac{1}{N}\sum_i|y_{i,t}-\hat{y}_{i,t}|$):
кризисные годы дают отчётливые выбросы ошибки у всех моделей; ансамбль
сохраняет наименьший MAE даже в кризисные периоды.

\begin{figure}[htbp]
  \centering
  \includegraphics[width=0.75\textwidth]{error_by_region}
  \caption{Топ-20 регионов по~MAE прогноза ансамбля CB$\times$0{,}25+LGB$\times$0{,}75.}
  \label{fig:error_region}
\end{figure}

\begin{figure}[htbp]
  \centering
  \includegraphics[width=0.95\textwidth]{error_heatmap}
  \caption{Тепловая карта ошибок прогноза (регион $\times$ год)
           для 30 субъектов с~наибольшим средним MAE.}
  \label{fig:error_heatmap}
\end{figure}

\subsection{Тест Дибольда--Мариано}

Для статистической верификации превосходства бустинговых моделей над наивным
прогнозом применяется тест Дибольда--Мариано (DM)~\cite{diebold1995}.
Нулевая гипотеза: обе модели имеют одинаковую точность прогноза.
Тест-статистика строится на дифференциале квадратичных потерь:
\begin{equation}
  d_t = e^2_{1,t} - e^2_{2,t}, \quad
  \mathrm{DM} = \frac{\bar{d}}{\widehat{\mathrm{SE}}(\bar{d})},
\end{equation}
где $e_{j,t}$~--- ошибки модели~$j$ в момент~$t$,
$\bar{d}$~--- среднее дифференциала по тестовым наблюдениям.

\begin{center}
\small
\begin{tabular}{lrrrr}
\toprule
Сравнение & $\bar{d}$ & DM & $p$ (двуст.) & $p$ (одност.) \\
\midrule
CatBoost vs Наивный  & $\approx 0$ & $0{,}046$ & $0{,}963$ & $0{,}482$ \\
LightGBM vs Наивный  & $< 0$       & $6{,}247$ & $<0{,}001$ & $<0{,}001$ \\
\bottomrule
\multicolumn{5}{l}{\scriptsize $p$ (одност.) --- вероятность того, что модель~2 не лучше модели~1.}
\end{tabular}
\end{center}

LightGBM статистически достоверно превосходит наивный прогноз на уровне
$\alpha = 0{,}001$: $\mathrm{DM} = 6{,}25$, $p < 0{,}001$.
Для CatBoost разница не значима ($p = 0{,}48$), несмотря на положительный
RMSE-выигрыш~--- высокая дисперсия между fold'ами (особенно 2020 и 2022--2023)
снижает мощность теста. Поскольку ансамбль включает LightGBM с долей~$0{,}75$,
статистически подтверждённое превосходство LightGBM транслируется на
ансамблевую модель.

\subsection{Метрики в рублях}

Наряду с метриками в логарифмическом пространстве предсказания переведены
в абсолютные значения ВРП (млрд руб.) посредством обратного преобразования:
\begin{equation}
  \widehat{\mathrm{GRP}}_{i,t} =
  \exp\!\bigl(\log\mathrm{GRP}_{i,t-1} + \hat{y}_{i,t}\bigr).
\end{equation}
Дополнительно вычисляется средняя абсолютная процентная ошибка:
\begin{equation}
  \mathrm{MAPE} = \frac{1}{n}\sum_{i,t}
  \left|\frac{\mathrm{GRP}_{i,t} - \widehat{\mathrm{GRP}}_{i,t}}
  {\mathrm{GRP}_{i,t}}\right| \times 100\%.
\end{equation}
Ансамбль CB$\times$0{,}25+LGB$\times$0{,}75 показывает наименьший MAPE среди
всех протестированных моделей, сохраняя лидерство как в логарифмическом,
так и в абсолютном масштабе измерения.

\begin{figure}[htbp]
  \centering
  \includegraphics[width=0.65\textwidth]{ruble_metrics_table}
  \caption{Метрики прогноза в абсолютном масштабе (млрд руб.):
           RMSE, MAE и MAPE по всем моделям.}
  \label{fig:ruble_metrics}
\end{figure}

\begin{figure}[htbp]
  \centering
  \includegraphics[width=0.75\textwidth]{mape_by_model}
  \caption{MAPE (\%) в разрезе моделей. Ансамбль CB$\times$0{,}25+LGB$\times$0{,}75
           достигает наименьшей средней процентной ошибки.}
  \label{fig:mape}
\end{figure}

%% ---------------------------------------------------------------
%%  5. ОБСУЖДЕНИЕ РЕЗУЛЬТАТОВ
%% ---------------------------------------------------------------
\section{Обсуждение результатов}

\subsection{Сравнительный анализ моделей}

\textit{Победа ансамбля и LightGBM.} Взвешенный ансамбль CatBoost+LightGBM
с оптимальными весами 0{,}25/0{,}75 показал наилучший RMSE~=~0{,}0382;
вклад ансамблирования над лучшей одиночной моделью (LightGBM) составляет
0{,}0006 (около~1{,}5\% от его RMSE), что отражает высокую корреляцию
предсказаний двух бустинговых алгоритмов.
LightGBM (RMSE~=~0{,}0388) обошёл CatBoost (RMSE~=~0{,}0408), несмотря
на традиционное преимущество CatBoost при работе с категориальными
признаками. Возможная причина: LightGBM лучше использует монотонные
монетарные признаки (ключевая ставка, USD/RUB) на непрерывной числовой
шкале, тогда как ordered boosting CatBoost ориентирован на нелинейные
паттерны в категориальных данных, которые здесь менее значимы.
Гиперпараметры обоих алгоритмов оптимизировались с помощью библиотеки Optuna. 60 итераций с внутренней валидацией через метод расширяющегося окна на данных за последние два года обучающей выборки.

\textit{Провал GPBoost.} GPBoost показал наихудшие результаты
($\mathrm{RMSE} = 0{,}488$, $R^2_{\mathrm{OOS}} = -280{,}9$) с
экстремально высокой дисперсией между фолдами (RMSE~2016~$\approx$~2{,}0).
Причина технического характера кроется в генерации прогнозов на тестовой
выборке гауссовская процессная компонента модели отключается
(\texttt{ignore\_gp\_model=True}), так как условные предсказания
случайных эффектов для новых временны\'{х} периодов неустойчивы.
Так или иначе, GPBoost деградирует до обычного градиентного бустинга без
каких-либо случайных эффектов, но также проигрывает LightGBM и CatBoost,
что используют в свою очередь ту же парадигму с более подходящими гиперпараметрами.

\textit{Линейные и нейросетевые базовые модели.} Ridge с FE
($\mathrm{RMSE}=0{,}052$, $R^2_{\mathrm{OOS}}=-0{,}62$) уступает
наивному прогнозу. Это объясняется тем, что within-трансформация удаляет существенную часть
межрегиональной дисперсии на коротком используемом панельном ряду (10--18 наблюдений
на регион), оставляя лишь~$\sim$10 внутрирегиональных отклонений для
регрессии с~$\sim$80 признаками. EmbMLP в свою очередь ($\mathrm{RMSE}=0{,}049$,
$R^2_{\mathrm{OOS}}=-0{,}44$) и вовсе чрезмерно страдает от классической проблемы малой
выборки. Для нее имеющиеся $85 \times 18 \approx 1\,500$ наблюдений попросту недостаточно. AR(1) ($\mathrm{RMSE}=0{,}049$) нестабилен
из-за малого числа точек в обучающем окне.

\subsection{Доминирование макропеременных}

SHAP-анализ обнаруживает, что наиболее значимыми предикторами темпа роста
ВРП являются монетарные и валютные индикаторы, а не отраслевые показатели
производства. Это экономически интерпретируемо: в условиях централизованной
денежно-кредитной политики России все регионы реагируют на федеральные
монетарные шоки практически однородно, тогда как отраслевая диверсификация
проявляется на более коротких временны\'{х} горизонтах, слабее заметных на
годовых данных. Кризисный индикатор 2020~г. в тройке лидеров~---
свидетельство того, что модель «научилась» значительному снижению темпа
роста в пандемийный год и корректно применяет этот эффект при наличии
dummy-признака.

Групповой SHAP показывает, что MACRO\_BLOCK и PRICE\_BLOCK доминируют
во всех пяти типах регионов (0{,}21--0{,}27 и 0{,}16--0{,}21
соответственно), подтверждая однородную реакцию региональных экономик
на федеральные монетарные шоки. Производственный блок сильнее для
сырьевых регионов (0{,}17 vs~0{,}08--0{,}11 у остальных) --- ожидаемо,
поскольку ВРП таких субъектов напрямую зависит от объёмов добычи.
Инвестиционный блок несколько выше у сервисных (0{,}16) и смешанных
(0{,}14) регионов. Переменная CATEGORICAL\_BLOCK имеет абсолютно нулевой вклад во всех
типах~--- это говорит нам о том, что принадлежности региона к тому или иному типу
не несёт никакой самостоятельной предсказательной силы после учёта
количественных признаков.

\subsection{Ограничения исследования}

\begin{enumerate}
  \item Чрезмерно малая панель, что
        ограничивает применение более сложных архитектур. Для каждого региона представлено всего около 20 наблюдего
  \item Часть показатели доступны исключительно с лагом публикации, что на самом деле завышает
        оценку практической применимости.
  \item Имеющиеся структурные разрывы (кризисы 2008, 2014, 2020 и 2022 годов) затрудняют
        перенос модели между экономическими режимами.
  \item Типология регионов задана статически, основываясь на исторической сводке, и никак не отражает эволюции
        экономической структуры.
  \item SHAP-значения объясняют прогноз модели, но ни коим образом не доказывают
        их причинности.
  \item GP-компонента GPBoost не используется при OOS-прогнозировании,
        что в свою очередь исключает ключевое преимущество
        данного алгоритма.
\end{enumerate}

%% ---------------------------------------------------------------
%%  6. ЗАКЛЮЧЕНИЕ
%% ---------------------------------------------------------------
\section{Заключение}

В рамках данной диссертации задача прогнозирования ВРП субъектов Российской Федерации
сформулирована как регрессия на панельных временных данных с горизонтом
$h=1$. Реализован воспроизводимый pipeline: предобработка без утечки
данных, скользящая кросс-валидация с расширяющимся окном (9~fold'ов,
тестовые годы 2015--2023) и сопоставление восьми моделей~--- наивного
прогноза, AR(1), Ridge с FE, EmbMLP, LightGBM, CatBoost, GPBoost
и ансамбля CatBoost+LightGBM.

Главный результат: \textbf{ансамбль CatBoost+LightGBM с весами 0{,}25/0{,}75
превзошёл все модели} по RMSE, MAE и $R^2_{\mathrm{OOS}}$.
RMSE~=~0{,}0382~--- снижение относительно наивного прогноза на 8{,}8\%
(против 0{,}0419); $R^2_{\mathrm{OOS}} = 0{,}144$ подтверждает
устойчивое превосходство над бенчмарком вне обучающей выборки.
Среди одиночных моделей LightGBM (RMSE~=~0{,}0388) незначительно
опережает CatBoost (RMSE~=~0{,}0408). Ablation study показал, что
добавление производственного блока к лаговым переменным даёт
наибольший прирост точности для LightGBM. Два robustness check~---
без структурных выбросов (Москва, ХМАО, ЯНАО, Сахалинская обл.)
и в разрезе кризисных периодов~--- подтвердили устойчивость
ранжирования бустинговых моделей.

SHAP-интерпретация LightGBM (лучшей одиночной модели) обнаружила
доминирование блока PRICE\_BLOCK: волатильность курса USD/RUB
(\texttt{usd\_rub\_volatility\_lag1}) оказалась признаком~№1 с вкладом,
вдвое превышающим следующий по важности предиктор. Макро-монетарные
признаки (MACRO\_BLOCK: ключевая ставка номинальная и реальная,
ставка на конец года) занимают позиции~3--4 и~6. CATEGORICAL\_BLOCK
(тип региона) имеет нулевой вклад во всех типах~--- предсказательная
сила полностью исходит из количественных характеристик, а не из
номинальной классификации.

GPBoost показал наихудшие результаты ($\mathrm{RMSE} = 0{,}488$,
$R^2_{\mathrm{OOS}} = -280{,}9$): техническое ограничение~---
отключение GP-компоненты при OOS-прогнозировании~--- де-факто
лишает модель её методологического преимущества.

Практическая значимость: прогностические системы для региональных
финансовых органов следует строить на ансамбле LightGBM+CatBoost
с полным набором признаков; ключевыми опережающими индикаторами
являются волатильность валютного курса и параметры денежно-кредитной
политики, а не отраслевые производственные индексы.

Статистическое превосходство LightGBM над наивным прогнозом подтверждено
тестом Дибольда--Мариано ($\mathrm{DM} = 6{,}25$, $p < 0{,}001$);
для ансамбля это превосходство наследуется от LGB-компоненты с долей~0{,}75.

Перспективные направления: исследование устойчивости при увеличении
горизонта прогноза ($h=2,3$), включение высокочастотных данных через
MIDAS-структуру, прямая верификация превосходства ансамбля над LightGBM
тестом Дибольда--Мариано с поправкой на множественное сравнение.

%% ---------------------------------------------------------------
%%  Список литературы
%% ---------------------------------------------------------------
\newpage
\begin{thebibliography}{20}

\bibitem{rosstat2024}
Федеральная служба государственной статистики. Валовой региональный
продукт за 2024 год. URL:
\url{https://rosstat.gov.ru/folder/313/document/279001}.

\bibitem{chi2025}
Chi T.\,C., Fan T.-H., Ghigliazza R.\,M., Giannone D., Wang~Z.
Macroeconomic Forecasting and Machine Learning.
\textit{arXiv:2510.11008}, 2025.

\bibitem{cerqua2025}
Cerqua A., Letta M., Pinto G.
On the (Mis)Use of Machine Learning With Panel Data.
\textit{Oxford Bulletin of Economics and Statistics}, 2025.

\bibitem{badykova2026}
Бадыкова И.\,Р. Валовой региональный продукт: анализ детерминант
для России с использованием моделей машинного обучения.
\textit{Экономика региона}, 2026. Т.~22. №~1. С.~97--108.
DOI: 10.17059/ekon.reg.2026-1-8.

\bibitem{lashina2023}
Lashina M., Grishunin S.
Comparison of Forecasting Power of Statistical Models for GDP Growth
Under Conditions of Permanent Crises.
\textit{Procedia Computer Science}, 2023. Vol.~221. P.~442--449.
DOI: 10.1016/j.procs.2023.07.059.

\bibitem{sigrist2022}
Sigrist F. Gaussian Process Boosting.
\textit{Journal of Machine Learning Research}, 2022.
Vol.~23. No.~232. P.~1--46.

\bibitem{smirnov2025}
Смирнов С.\,В. Методы машинного обучения в макроэкономическом
прогнозировании: предварительные итоги.
\textit{Вопросы экономики}, 2025. №~10. С.~131--154.
DOI: 10.32609/0042-8736-2025-10-131-154.

\bibitem{tsukarev2025}
Цукарев Т., Погосян К., Лемба К., Гришин Д.
Наукастинг ВВП: от традиционных эконометрических моделей к машинному
обучению. \textit{Рабочий документ ЕФСР}, РД/25/1, 2025.
URL: \url{https://efsd.org/research/working-papers/}.

\bibitem{makeeva2022}
Макеева Н.\,М., Станкевич И.\,П.
Наукастинг элементов использования ВВП России.
\textit{Экономический журнал ВШЭ}, 2022. Т.~26. №~4. С.~598--622.
DOI: 10.17323/1813-8691-2022-26-4-598-622.

\bibitem{malkina2025}
Малкина М.\,Ю., Сочков А.\,Л., Капустина Ю.\,И.
Прогнозирование ВРП региона с применением методов искусственного
интеллекта.
\textit{Journal of New Economy}, 2025. Т.~26. №~3. С.~67--85.
DOI: 10.29141/2658-5081-2025-26-3-4.

\bibitem{semiturkin2022}
Semiturkin O., Shevelev A.
Forecasting Regional Inflation Rates Using Machine Learning Methods:
The Case of Siberia Macroregion.
\textit{Bank of Russia Working Paper Series}, No.~91, 2022.

\bibitem{chernyshova2022}
Chernyshova G.\,Yu., Rasskazkin N.\,D.
Software Implementation of Ensemble Models for the Analysis of
Regional Socio-Economic Development Indicators.
\textit{Izvestiya of Saratov University. Mathematics. Mechanics.
Informatics}, 2022. Vol.~22. Iss.~1. P.~130--137.
DOI: 10.18500/1816-9791-2022-22-1-130-137.

\bibitem{buckmann2025}
Buckmann M., Potjagailo G.
Infusing Economically Motivated Structure into Machine Learning Methods.
\textit{Bank of England Staff Working Paper}, 2025.

\bibitem{attolico2025}
Attolico L.
Explainable Machine Learning for Macroeconomic and Financial
Nowcasting. \textit{arXiv:2512.00399}, 2025.

\bibitem{lundberg2017}
Lundberg S.\,M., Lee S.-I.
A Unified Approach to Interpreting Model Predictions.
\textit{Advances in Neural Information Processing Systems}, 2017.

\bibitem{breiman2001}
Breiman L. Random Forests.
\textit{Machine Learning}, 2001. Vol.~45. P.~5--32.

\bibitem{ke2017}
Ke G., Meng Q., Finley T., et al.
LightGBM: A Highly Efficient Gradient Boosting Decision Tree.
\textit{Advances in Neural Information Processing Systems}, 2017.

\bibitem{dorogush2018}
Dorogush A.\,V., Ershov V., Gulin A.
CatBoost: Gradient Boosting with Categorical Features Support.
\textit{arXiv:1810.11363}, 2018.

\bibitem{diebold1995}
Diebold F.\,X., Mariano R.\,S.
Comparing Predictive Accuracy.
\textit{Journal of Business \& Economic Statistics}, 1995.
Vol.~13. No.~3. P.~253--263.
DOI: 10.1080/07350015.1995.10524599.

\end{thebibliography}

\end{document}
